\documentclass[11pt,a4paper]{article}

\usepackage[brazil]{babel}
\usepackage[T1]{fontenc}
\usepackage{lmodern}
\usepackage{microtype}
\usepackage[margin=2.2cm]{geometry}
\usepackage{amsmath}
\usepackage{booktabs}
\usepackage{tabularx}
\usepackage{enumitem}
\usepackage{tikz}
\usepackage{wrapfig}
\usepackage{algorithm}
\usepackage{algpseudocode}
\usepackage[hidelinks]{hyperref}
\usepackage{xcolor}

\definecolor{destaque}{HTML}{1F4E79}
\setlist{nosep, leftmargin=*}
\setlength{\parskip}{0.4em}
\setlength{\parindent}{0pt}

% Frase de abertura de cada seção: conecta com a anterior.
\newcommand{\gancho}[1]{\textcolor{destaque}{\textit{#1}}\par}

\floatname{algorithm}{Algoritmo}
\algrenewcommand\algorithmicrequire{\textbf{Entrada:}}
\algrenewcommand\algorithmicensure{\textbf{Saída:}}
\algrenewcommand\algorithmicfor{\textbf{para cada}}
\algrenewcommand\algorithmicdo{\textbf{faça}}
\algrenewcommand\algorithmicif{\textbf{se}}
\algrenewcommand\algorithmicthen{\textbf{então}}
\algrenewcommand\algorithmicend{\textbf{fim}}

\title{\textbf{Ativação dinâmica de áreas por grafo}\\[0.3em]
  \large Justificativa do algoritmo}
\author{Caio Trindade e Luisa Ferreira\\
  \small Inteligência Artificial para Jogos I --- Representação do Mundo}
\date{Setembro de 2026}

\begin{document}
\maketitle

\section{O problema}

\gancho{Por que não dá para simplesmente atualizar tudo de uma vez?}

O nível tem 9 áreas povoadas por inimigos, e os testes de balanceamento pedem
chegar a milhares deles. Atualizar todos a cada quadro custa $O(N)$, com $N$ o
total de inimigos, e $N$ cresce com o tamanho do mundo. Mas o jogador só
interage com o que está perto: atualizar um inimigo do outro lado do mapa é
trabalho que ninguém observa. É a pergunta que a aula faz sobre malhas
\emph{"o que é preciso atualizar a cada iteração do jogo?"} e a resposta é:
só a vizinhança do jogador.

\section{A modelagem: o mundo como grafo}

\gancho{Vizinhança é uma relação entre lugares; a estrutura natural para
relações é um grafo.}

\begin{wrapfigure}{r}{0.30\textwidth}
  \vspace{-1.2em}
  \centering
  \begin{tikzpicture}[
      area/.style={circle, draw=destaque, thick, minimum size=6mm,
                   inner sep=0pt, font=\small},
      scale=0.95]
    \foreach \i in {0,...,8} {
      \pgfmathtruncatemacro{\x}{mod(\i,3)}
      \pgfmathtruncatemacro{\y}{div(\i,3)}
      \node[area] (n\i) at (\x*1.3,\y*1.3) {\i};
    }
    \foreach \a/\b in {0/1,1/2,3/4,4/5,6/7,7/8,0/3,3/6,1/4,4/7,2/5,5/8}
      \draw[thick] (n\a) -- (n\b);
    \foreach \a/\b in {0/4,1/5,3/7,4/8,1/3,2/4,4/6,5/7}
      \draw[thin, gray] (n\a) -- (n\b);
  \end{tikzpicture}
  \par\vspace{0.3em}{\small Malha $3\times3$ com diagonais (cinza): grau 3 nos cantos, 5 nas bordas, 8 no centro.\par}
  \vspace{-1.5em}
\end{wrapfigure}

O mundo é um grafo não dirigido $G = (V, A)$. Cada \textbf{vértice} é uma área:
uma região retangular com o próprio conteúdo (inimigos, itens) e o próprio
estado (ativa ou inativa). Cada \textbf{aresta} liga duas áreas que se tocam,
por um lado ou por um canto: são as áreas entre as quais se pode passar
diretamente. É a subdivisão do mundo vista como grafo, como na
aula: cada região vira um nó, e regiões vizinhas são ligadas.

\textbf{Representação: lista de adjacência.} A única operação que o algoritmo
faz sobre o grafo é \emph{visitar os vizinhos de um vértice}, e é exatamente
essa a operação que a lista de adjacência torna eficiente: custa $O(d)$, com
$d$ o grau do vértice. Numa matriz de adjacência a mesma consulta varre uma
linha inteira, $O(|V|)$, e a matriz ocupa $O(|V|^2)$ para um grafo esparso.

\textbf{Por que grafo e não só a malha.} Na malha, os vizinhos são
\emph{calculados} a partir da posição; no grafo, são um \emph{dado
armazenado}. Isso desacopla a topologia da geometria: uma parede entre duas
áreas vizinhas, um corredor ou portal ligando áreas distantes, áreas de tamanhos
diferentes. A malha regular vira um caso particular do grafo. Hoje o grafo é
gerado como malha, mas mudar a topologia significa mudar só a função que o
constrói; o algoritmo e a análise abaixo continuam valendo.

\section{O algoritmo: a janela ativa}

\gancho{Com o mundo como grafo, o que o jogo faz a cada quadro?}

O estado mantido é o conjunto de \textbf{áreas ativas}, com no máximo $K = 4$
elementos. A cada quadro ele é recalculado a partir da posição do jogador:

\begin{algorithm}[h]
\caption{Atualização das áreas ativas}
\begin{algorithmic}[1]
\Require posição $P$ do jogador, grafo $G$, limiar $D$, limite $K$
\State $a \gets$ área que contém $P$ \Comment{se não houver, nada muda}
\State $C \gets \{a\}$
\For{vizinho $v$ de $a$ em $G$} \Comment{percurso local: profundidade 1}
  \If{$\operatorname{dist}(P, \text{limites}(v)) < D$} $C \gets C \cup \{v\}$ \EndIf
\EndFor
\If{$|C| > K$} manter as $K$ mais próximas de $P$, sempre com $a$ \EndIf
\State desativar as áreas que saíram de $C$; carregar e ativar as que entraram
\end{algorithmic}
\end{algorithm}

Depois disso, só os inimigos das áreas ativas são atualizados (perseguição e
dano); as demais ficam paradas. O passo central é a linha 3: uma \textbf{busca
em largura limitada à profundidade 1} a partir da área atual. Ela visita o nó e
seus vizinhos imediatos e para, sem nunca percorrer o grafo inteiro.

\section{O custo}

\gancho{Por que isso escala com o tamanho do mundo?}

Sejam $n$ o número de áreas, $d$ o grau máximo (8 na malha), $N$ o total de
inimigos e $K$ o limite de áreas ativas.

\begin{center}
\begin{tabular}{@{}lll@{}}
\toprule
Etapa & Ingênuo & Janela ativa \\
\midrule
Decidir o que atualizar & --- & $O(d \log d)$, constante em $n$ \\
Atualizar inimigos & $O(N)$ & $O(K \cdot N/n)$ \\
\bottomrule
\end{tabular}
\end{center}

A simulação cai para a fração $K/n$ do total: $4/9 \approx 44\%$ na malha
$3\times3$, $4\%$ numa $10\times10$, $0{,}4\%$ numa $32\times32$. Na $3\times3$
o ganho é modesto (cerca de $2{,}2\times$), porque quatro áreas já são quase
metade do mundo; \textbf{o ganho é assintótico}, cresce com a malha. E $K$
cheio é o pior caso: longe das divisas só a área atual fica ativa.

\section{Correção e limites}

\gancho{Barato só vale se estiver certo. O que precisa valer ao fim de todo
quadro?}

\begin{enumerate}
  \item A área que contém o jogador está sempre ativa.
  \item Há no máximo $K$ áreas ativas.
  \item Toda área ativa é a atual ou vizinha dela no grafo.
  \item Nenhuma área ativa está descarregada da memória.
  \item Nenhum inimigo de área inativa se moveu.
\end{enumerate}

O invariante 3 é o que liga o algoritmo ao grafo e garante que o custo é local.
É também o que exige as \textbf{arestas diagonais}: perto de um canto, o
jogador está perto de três áreas além da sua, e uma delas só toca a área atual
pelo canto. Sem a aresta diagonal ela ficaria a dois saltos, fora do alcance da
profundidade 1, e o máximo de 4 áreas ativas do enunciado --- que é exatamente
esse caso do canto --- nunca aconteceria. A condição que resta é: \textbf{o limiar $D$ deve ser menor que a menor
dimensão de uma área}. Assim, para chegar perto de uma área a duas arestas de
distância, o jogador precisa atravessar a do meio, que nesse momento já é a
atual; a profundidade 1 é exata, não uma aproximação. Se $D$ passar disso, o
percurso precisa virar uma busca em largura completa, com fila e visitados,
cortada por distância.

O limiar também evita \textbf{oscilação na fronteira}: se a ativação dependesse
só de ``estar dentro da área'', andar sobre uma divisa ligaria e desligaria a
vizinha a cada quadro. Com $D$, a área vizinha acorda antes de o jogador chegar.
Pelo mesmo motivo, liberar uma área da memória precisa de um atraso: a
implementação usa 30 segundos inativa antes de descarregar (\texttt{areaUnloadDelay}),
e o descarte é destrutivo --- inimigos mortos e itens coletados voltam ao
estado inicial se a área for recarregada, já que nenhum resumo do progresso é
guardado.

\section{Alternativas descartadas}

\gancho{Por que não outra estrutura vista em aula?}

\begin{center}
\begin{tabularx}{\textwidth}{@{}lX@{}}
\toprule
Abordagem & Por que não \\
\midrule
Atualizar tudo & $O(N)$ por quadro, justamente o que se quer evitar. \\
Distância a cada inimigo & Ainda $O(N)$: medir a distância custa o mesmo que atualizar. \\
Quadtree & Responde bem \emph{o que está perto} (consulta geométrica). O problema
  pergunta \emph{o que está conectado}; as duas coincidem num mundo aberto, mas
  deixam de coincidir com paredes ou passagens. \\
Matriz de adjacência & Listar vizinhos custa $O(n)$ e o espaço é $O(n^2)$ para
  um grafo esparso. \\
Busca em largura completa & Correta, mas desnecessária enquanto $D$ for menor
  que a área: profundidade 1 basta. \\
\bottomrule
\end{tabularx}
\end{center}

\section{Como demonstrar}

\gancho{O argumento fica completo com números.}

Medido o tempo médio de quadro (média de 60 quadros simulados, sem terminal
nem espera de taxa de quadros) variando o número de inimigos, em duas
condições: com a janela ativa e com todas as áreas forçadas a ativas (a linha
de base ingênua, $O(N)$). O jogador foi posicionado no pior caso de $K$ cheio:
o ponto interno onde quatro áreas se tocam. Os números brutos, incluindo mais
pontos que os da tabela abaixo, estão em \texttt{WorldGraphGame/benchmarks/frame\_time.csv}.

\begin{center}
\begin{tabular}{@{}lrrrr@{}}
\toprule
Malha & $N$ (inimigos) & Ingênuo (ms) & Janela ativa (ms) & Ganho \\
\midrule
$3\times3$  & 900    & 0,0104 & 0,0044 & $2{,}4\times$ \\
$3\times3$  & 4500   & 0,0334 & 0,0130 & $2{,}6\times$ \\
$3\times3$  & 9000   & 0,0518 & 0,0337 & $1{,}5\times$ \\
$10\times10$ & 10000  & 0,0524 & 0,0017 & $30{,}8\times$ \\
$10\times10$ & 50000  & 0,2705 & 0,0074 & $36{,}6\times$ \\
$10\times10$ & 100000 & 0,7874 & 0,0224 & $35{,}2\times$ \\
\bottomrule
\end{tabular}
\end{center}

O resultado confirma a previsão da Seção 4: a linha de base cresce
linearmente com $N$, a janela ativa cresce muito mais devagar, e o ganho na
malha $3\times3$ fica na faixa do $2{,}2\times$ estimado a partir de $K/n =
4/9$ (o ponto de $9000$ destoa por ruído de medição --- o processo compete
por CPU com o resto do sistema, e o efeito pesa mais nos tempos já pequenos
da janela ativa). Na malha $10\times10$ o ganho medido (entre $31\times$ e
$37\times$) fica até acima do $25\times$ estimado a partir de $K/n =
4/100$, o que é consistente: $K/n$ é o pior caso, e o ganho real tende a
ficar igual ou melhor que ele. O comportamento assintótico --- janela ativa
crescendo muito mais devagar que a linha de base --- é o que a Seção 4
previu.

\end{document}
